\begin{afterword}
\summaryhead

The post opens with three large projects that ran late and over budget: Denver International Airport, the Eurofighter Typhoon and the Sydney Opera House. It grants that such cases may reflect selective attention or bad incentives, but says that experiments with individual planners show the same bias. Students asked for dates by which they were 50\%, 75\% and 99\% sure to finish a project met those dates only 13\%, 19\% and 45\% of the time. Another study found that people's ``realistic'' scenarios were the same as their best-case ones.

The remedy is the ``outside view'': ignore the special features of the task and ask how long broadly similar tasks have taken, or ask an experienced outsider. Christmas shoppers who planned in detail were more optimistic than those who did not, and Japanese students' past record predicted their essay times better than their plans did. The post warns that the outside answer will sound too long, and says that it is true.

\responsehead

The post reports real research, and its main claim holds up in the studies it cites: people's forecasts of their own finishing times lean toward optimism, and looking at past cases can remove the lean. The problems are at the two ends of the post, and in what it leaves out of its own sources.

The opening does not bear on the claim. Three famous overruns take the first three paragraphs. The post then grants that they are ``very probably'' explained by selective attention or incentives, and later says its fix will not work on projects of that size. For Denver, the government's auditors put most of the cost growth down to changes in the airport's scope. The cases introduce the topic; the evidence for a bias in individual planners is the student studies.

The worst-case claim is asserted where it could have been shown. The post says, with no source, that reality usually turns out somewhat worse than the worst case, and later says ``we saw'' that picturing the worst case is not enough. The 99\% forecasts come close to showing it. Direct evidence was in the post's own sources: the paper in its fifth footnote reports that pessimistic scenarios did not change people's predictions, and the 1994 paper in its second footnote found that about half of one class of students finished after their own worst-case dates. The post reports neither finding.

The ending claims more than the evidence gives. The 1994 paper tested versions of both remedies the post recommends. Students who only recalled their past assignments stayed optimistic. Students asked to relate their past to the task lost the optimistic lean, but their forecasts were no more accurate. Fellow students who forecast for others were more pessimistic, and also no more accurate. That supports ``a fairly reliable way'' to correct optimism. It does not support ``much more accurate'' outsiders, or ``This answer is true.'' An outside estimate corrects the average; it can still miss the case at hand, in either direction.

In short: a fair summary of real findings on optimistic forecasts, opened with megaprojects the post itself sets aside, and closed with promises (``much more accurate'' outsiders, an answer that ``is true'') that the one cited source I could check, the 1994 paper, did not find.
\end{afterword}
